\documentclass[11pt]{amsart}

\usepackage[margin=1.15in]{geometry}
\usepackage{amsmath,amssymb,amsthm,mathtools}
\usepackage{microtype}
\usepackage{hyperref}
\usepackage[nameinlink,capitalize]{cleveref}

\hypersetup{
  colorlinks=true,
  linkcolor=blue,
  citecolor=blue,
  urlcolor=blue
}

\theoremstyle{plain}
\newtheorem{theorem}{Theorem}[section]
\newtheorem{proposition}[theorem]{Proposition}
\newtheorem{lemma}[theorem]{Lemma}
\newtheorem{corollary}[theorem]{Corollary}

\theoremstyle{definition}
\newtheorem{definition}[theorem]{Definition}
\newtheorem{hypothesis}[theorem]{Hypothesis}

\theoremstyle{remark}
\newtheorem{remark}[theorem]{Remark}

\numberwithin{equation}{section}

\title[Seregin limits and Type I exclusion]{A Residual-class Criterion for Excluding Type I Navier--Stokes Singularities}
\author{Guillem Duran Ballester}
\date{}

\subjclass[2020]{35Q30, 35B44, 35B40, 76D05}
\keywords{Navier--Stokes equations, Type I singularities, ancient solutions, Seregin limits, Liouville theorems}

\begin{document}

\begin{abstract}
We give a conditional reduction for Type I singularities of finite-energy
suitable weak solutions of the three-dimensional Navier--Stokes equations.
Starting from the positive pressure-normalized Caffarelli--Kohn--Nirenberg
concentration carried by every singular point, we extract ancient limits by
the Seregin blow-up procedure and separate the possible centered limits into
small-amplitude, stationary \(L^3\), uniformly \(L^3\)-tight, fast-decaying,
structured, and residual alternatives.  The small-amplitude and stationary
\(L^3\) alternatives are excluded unconditionally by classical Liouville
theorems.  The structured taxonomy also includes the established
Lei--Zhang--Zhao \(L_t^\infty L_x^p\)-swirl branch and the
Lei--Ren--Zhang periodic bounded-swirl branch.  Assuming Liouville criteria
for the remaining tight and structured alternatives, full finite-energy Type I
exclusion is reduced to one explicit analytic assertion: the residual
non-tight/non-decaying class is empty, or else satisfies a replacement
Liouville theorem.  The positive concentration and the local Type I/Type II
dichotomy are invoked as results proved earlier in the series; the additional
inputs are precisely the stated nonclassical Liouville criteria.
\end{abstract}

\maketitle

\section{Introduction}

Let \((u,p)\) be a finite-energy suitable weak solution of the
three-dimensional incompressible Navier--Stokes equations on
\(\mathbb R^3\times(0,T)\),
\begin{equation}\label{eq:ns}
   \partial_tu+u\cdot\nabla u+\nabla p=\Delta u,\qquad
   \nabla\cdot u=0 .
\end{equation}
A finite-time singular point \(z_*=(x_*,T)\) is called Type I, in the sense
used here, if the scale-invariant \(L^\infty\) quantity remains bounded near
\(T\):
\begin{equation}\label{eq:type-i}
   \limsup_{t\uparrow T}\sqrt{T-t}\,
   \|u(\cdot,t)\|_{L^\infty(\mathbb R^3)}<\infty .
\end{equation}
Equivalently, after decreasing the time interval if necessary, there are
constants \(M<\infty\) and \(t_0<T\) such that
\[
   \|u(\cdot,t)\|_{L^\infty(\mathbb R^3)}
   \le M(T-t)^{-1/2}
   \qquad\hbox{for }t_0<t<T.
\]
This is the form used when passing the Type I bound to blow-up sequences.
Indeed, if the limsup in \eqref{eq:type-i} equals \(L<\infty\), choose
\(M>L\); by the definition of limsup there is \(t_0<T\) such that
\(\sqrt{T-t}\,\|u(t)\|_{L^\infty}\le M\) for \(t_0<t<T\).
The local Caffarelli--Kohn--Nirenberg concentration theorem gives a sequence
of shrinking backward parabolic cylinders centered at every singular point on
which the pressure-normalized critical density is bounded from below.  In the
Type I case, the Seregin blow-up procedure applied to such a concentration
sequence produces a nonzero bounded ancient solution in backward
self-similar variables,
\begin{equation}\label{eq:centered}
   \partial_\tau V+(V\cdot\nabla)V+\nabla P
   =
   \Delta V-\frac12(V+y\cdot\nabla V),
   \qquad \nabla\cdot V=0,
\end{equation}
defined on \(\mathbb R^3\times\mathbb R\).  One first extracts an
unrenormalized ancient suitable weak limit and then passes to centered
variables,
\[
   V(y,\tau)=\sqrt{-t}\,U(y\sqrt{-t},t),\qquad \tau=-\log(-t).
\]
The inherited positive local critical density prevents the centered limit from
vanishing identically.

\begin{remark}[Source-level critical \(L^3\) regime is already closed]
\label{rem:source-L3-closed}
If the original finite-energy suitable weak solution satisfies
\[
   \sup_{0<t<T}\|u(t)\|_{L^3(\mathbb R^3)}<\infty ,
\]
then the endpoint theorem of Escauriaza--Seregin--Sverak \cite{ESS2003}
implies that no point of \(\mathbb R^3\times\{T\}\) is singular.  Thus a
Type I singularity would have to lie outside the global critical
\(L^\infty_tL^3_x\) source regime.  This is a condition on the original
solution, not on an individual ancient limit, and it is therefore not included
among the residual classes below.
\end{remark}

No unconditional Liouville theorem is known for all bounded ancient solutions
of \eqref{eq:centered}.  The reduction below is therefore organized around
classes for which Liouville mechanisms are available or are formulated below
as explicit hypotheses.  Very small bounded ancient solutions vanish by the mild
formulation and decay of the centered Stokes semigroup.  Stationary
self-similar profiles in \(L^3(\mathbb R^3)\) vanish by the theorem of
Necas--Ruzicka--Sverak.  Uniform \(L^3\)-tightness is the natural compactness
condition for minimal-element and compact-hull arguments.  Fast decay and
geometric structure give additional routes to rigidity.  What remains is the
residual non-tight/non-decaying class.  The main theorem proves that excluding
this residual class, together with the stated Liouville inputs for the
nonclassical classes, excludes finite-energy Type I singularities.

\subsection{Logical ledger of the reduction}

The argument below is a finite exhaustion argument.  It relies on the following
inputs.  First, local CKN concentration at a singular point
supplies a positive pressure-normalized density.  Second, the Type I bound and
Seregin compactness yield a nonzero bounded
ancient solution of \eqref{eq:centered}.  Third, the source-level
\(L^\infty_tL^3_x\) regime has already been removed by
Escauriaza--Seregin--Sverak.  Fourth, small ancient limits and stationary
\(L^3\) ancient limits are unconditionally zero.  Fifth, the tight,
structured, and residual classes are eliminated precisely under the hypotheses
explicitly stated below.  Thus the final theorem is conditional precisely in
those places where a nonclassical Liouville theorem is still required.

\section{Local concentration and Seregin ancient limits}

\begin{definition}[Finite-energy Type I source]\label{def:source}
A pair \((u,p)\) is a finite-energy Type I source at \(z_*=(x_*,T)\) if
\((u,p)\) is a suitable weak solution of \eqref{eq:ns} on
\(\mathbb R^3\times(0,T)\) in the Caffarelli--Kohn--Nirenberg sense,
\[
   u\in L^\infty(0,T;L^2(\mathbb R^3))
      \cap L^2(0,T;\dot H^1(\mathbb R^3)),
   \qquad
   p\in L^{3/2}_{\mathrm{loc}}(\mathbb R^3\times(0,T)),
\]
the local energy inequality holds, \(z_*\) is a singular point, and the Type I
bound \eqref{eq:type-i} holds.
When the singular point is part of the notation, we write
\((u,p,z_*)\) for the corresponding source triple.
\end{definition}

Here a point is regular if the velocity is essentially bounded in some
backward parabolic cylinder with that point as top center; otherwise it is
singular.  This is the regularity notion used in the CKN epsilon-regularity
criterion.

For \(z_0=(x_0,t_0)\) and \(r>0\), write
\[
   Q_r(z_0)=B_r(x_0)\times(t_0-r^2,t_0).
\]
For an integrable function \(f\) on a ball \(B\), write
\[
   (f)_B=\frac1{|B|}\int_B f(x)\,dx .
\]
Define the pressure-normalized CKN quantities
\[
   C(z_0,r)=r^{-2}\int_{Q_r(z_0)}|u|^3\,dx\,dt
\]
and
\[
   D(z_0,r)=r^{-2}\int_{t_0-r^2}^{t_0}
      \int_{B_r(x_0)}
      |p(x,t)-(p)_{B_r(x_0)}(t)|^{3/2}\,dx\,dt.
\]

Throughout the paper, pressures are understood modulo functions of time.  Thus
two pressures \(p\) and \(p-b(t)\) represent the same velocity equation and the
same local energy inequality.  All pressure norms used below are either
pressure oscillations over balls or are written after subtracting an explicitly
allowed time-dependent gauge.  This convention is the usual CKN pressure
normalization and is justified for the rescaled solutions by
\cite{Paper1TypeI}, lemma \texttt{lem:gauges-\allowbreak harmless}.

\begin{theorem}[Local concentration at singular points]\label{thm:local-concentration}
Let \((u,p)\) be suitable in a parabolic neighborhood of
\(z_*=(x_*,T)\).  If \(z_*\) is singular, then there are \(\eta>0\) and
radii \(r_k\downarrow0\), with \(Q_{r_k}(z_*)\) contained in the neighborhood,
such that
\begin{equation}\label{eq:positive-ckn}
   C((x_*,T),r_k)+D((x_*,T),r_k)\ge\eta
   \qquad\hbox{for all }k.
\end{equation}
Moreover, along every such positive local concentration sequence, either a
local Type I scale-invariant \(L^\infty\) bound holds near \((x_*,T)\), or the
sequence lies in the local Type II alternative.
\end{theorem}

\begin{proof}
This is the local concentration entry result from
\cite{Paper0LocalConcentration}.  More precisely, the pressure-normalized
CKN criterion \cite{CKN1982} is recorded there as theorem
\texttt{thm:ckn}; theorem \texttt{thm:singular-\allowbreak positive} proves that every
singular point carries a positive concentration sequence in the sense of
\eqref{eq:positive-ckn}; and theorem \texttt{thm:blowup-\allowbreak rate-\allowbreak dichotomy} proves
the stated local Type I/Type II alternative along such shrinking cylinders.
The scale-invariance of the quantities \(C\) and \(D\), needed to identify the
same sequence after rescaling, is proposition \texttt{prop:critical-\allowbreak scaling}
in the same paper.  No global compactness or finite-energy input is used in
this local step.
\end{proof}

\begin{lemma}[Scaling and pressure gauges]\label{lem:scaling-gauges}
Let \((u,p)\) be suitable near \(z_*=(x_*,T)\), and set
\[
   u^{(\lambda)}(x,t)=\lambda u(x_*+\lambda x,T+\lambda^2t),
   \qquad
   p^{(\lambda)}(x,t)=\lambda^2 p(x_*+\lambda x,T+\lambda^2t).
\]
Then \((u^{(\lambda)},p^{(\lambda)})\) is suitable on the rescaled domain.
Moreover
\[
   C_{u^{(\lambda)}}((0,0),r/\lambda)=C_u(z_*,r),
   \qquad
   D_{p^{(\lambda)}}((0,0),r/\lambda)=D_p(z_*,r),
\]
where the pressure averages are taken over the corresponding balls.  Replacing
\(p^{(\lambda)}\) by \(p^{(\lambda)}-b(t)\), with \(b\) depending only on
time, does not change the equation, the local energy inequality, or the
quantity \(D\).
\end{lemma}

\begin{proof}
The suitability and scale-invariance assertions are proposition
\texttt{prop:critical-\allowbreak scaling} of \cite{Paper0LocalConcentration} and lemma
\texttt{lem:scaling} of \cite{Paper1TypeI}.  The pressure-gauge assertion is
lemma \texttt{lem:gauges-\allowbreak harmless} of \cite{Paper1TypeI}: subtracting a
function of time has zero spatial gradient, leaves the distributional equation
unchanged, and cancels from every pressure oscillation
\(p-(p)_{B_r}(t)\).  The cited proofs use the Navier--Stokes parabolic
scaling, the local energy inequality tested against pulled-back nonnegative
cutoffs, and the corresponding parabolic cylinder Jacobian.
\end{proof}

\begin{proposition}[Extraction of a normalized Seregin limit]
\label{prop:seregin-extraction}
Let \((u,p,z_*)\) be a finite-energy Type I source.  Then some positive
concentration sequence \(\lambda_k\downarrow0\) produces, after passing to a
subsequence and fixing pressure gauges, an ancient suitable weak limit
\((U,\Pi)\) on \(\mathbb R^3\times(-\infty,0)\).  The limit satisfies
\[
   \sup_{t<0}\sqrt{-t}\,\|U(\cdot,t)\|_{L^\infty(\mathbb R^3)}<\infty .
\]
Its centered representative \(V(y,\tau)=\sqrt{-t}\,U(y\sqrt{-t},t)\),
\(\tau=-\log(-t)\), is a Seregin ancient limit satisfying
\eqref{eq:nonvanishing}.
\end{proposition}

\begin{proof}
Choose \(\lambda_k\downarrow0\) from \Cref{thm:local-concentration}.  The
rescaled domains exhaust \(\mathbb R^3\times(-\infty,0)\), and
\Cref{lem:scaling-gauges} preserves suitability, the CKN normalization, and
pressure gauges.  The extraction itself is the classical Seregin
blow-up construction \cite{Seregin2012}, in the local Type I formulation of
Albritton--Barker \cite{AlbrittonBarker2019}.

All estimates used by that construction in the present normalization have
already been proved earlier in the series.  In \cite{Paper1TypeI}, lemma
\texttt{lem:inherited-bound} gives the compact-cylinder Type I bound for the
rescalings; lemma \texttt{lem:energy-diss} gives the local energy and
dissipation bounds; lemma \texttt{lem:pressure} fixes admissible pressure
gauges and gives local \(L^{3/2}\) pressure compactness; lemma
\texttt{lem:time-derivative} gives the negative-Sobolev time-derivative bound;
proposition \texttt{prop:compactness} gives strong local velocity convergence
and weak local pressure convergence; proposition \texttt{prop:ancient-limit}
passes the equation, divergence-free condition, and local energy inequality to
the ancient suitable weak limit; and lemmas \texttt{lem:limit-bound} and
\texttt{lem:smoothness} give the inherited Type I bound and smoothness.

The point specific to the present formulation is the normalization used below.
The positive local CKN density comes from \cite{Paper0LocalConcentration},
theorem \texttt{thm:singular-\allowbreak positive}.  Its persistence under the above
compactness and pressure gauges is supplied by \cite{Paper1TypeI}, lemma
\texttt{lem:scale-selection} and proposition \texttt{prop:nonzero}.  Passing
from \((U,\Pi)\) to centered variables is \cite{Paper1TypeI}, lemmas
\texttt{lem:renormalized-equation} and \texttt{lem:renorm-bounds}.  Under the
change of variables \((x,t)=(e^{-\tau/2}y,-e^{-\tau})\), the compact
negative-time cylinder carrying positive mass maps to
\(K\Subset\mathbb R^3\times\mathbb R\), and the pressure gauge becomes a
function \(a(\tau)\).  This is precisely \eqref{eq:nonvanishing}.
\end{proof}

\begin{definition}[Seregin ancient limit]\label{def:seregin-limit}
A smooth bounded ancient solution \(V\) of \eqref{eq:centered}, with pressure
understood modulo functions of \(\tau\), is called a Seregin ancient limit if
there exist a finite-energy Type I source \((u,p,z_*)\), scales
\(\lambda_k\downarrow0\) forming a positive local concentration sequence in
the sense of \eqref{eq:positive-ckn}, and the corresponding rescalings
\[
   u^{(k)}(x,t)=\lambda_k u(x_*+\lambda_k x,T+\lambda_k^2t),
   \qquad
   p^{(k)}(x,t)=\lambda_k^2 p(x_*+\lambda_k x,T+\lambda_k^2t),
\]
which converge locally, after passing to a subsequence and choosing pressure
gauges, to an ancient suitable weak solution \((U,\Pi)\) on
\(\mathbb R^3\times(-\infty,0)\).  The compactness is taken along a
concentration subsequence for which the positive pressure-normalized CKN
density persists in the limit.  The limit \(U\) satisfies
\[
   \sup_{t<0}\sqrt{-t}\,\|U(\cdot,t)\|_{L^\infty(\mathbb R^3)}<\infty ,
\]
and \(V\) is the self-similar representative of \(U\):
\[
   V(y,\tau)=\sqrt{-t}\,U(y\sqrt{-t},t),\qquad \tau=-\log(-t).
\]
The associated centered pressure class is represented by
\[
   P(y,\tau)=(-t)\,\Pi(y\sqrt{-t},t),
\]
and replacing \(\Pi\) by \(\Pi-b(t)\) replaces \(P\) by
\(P-e^{-\tau}b(-e^{-\tau})\).  Thus the pressure remains defined modulo a
function of \(\tau\), exactly as allowed by \cite{Paper1TypeI}, lemma
\texttt{lem:gauges-\allowbreak harmless}.
In particular, after passing to centered variables, there are a compact
cylinder \(K\Subset\mathbb R^3\times\mathbb R\), a pressure gauge \(a(\tau)\),
and \(\eta_0>0\) such that
\begin{equation}\label{eq:nonvanishing}
   \int_K\left(|V(y,\tau)|^3+|P(y,\tau)-a(\tau)|^{3/2}\right)\,dy\,d\tau
   \ge\eta_0 .
\end{equation}
\end{definition}

\begin{definition}[Normalized Seregin collection]\label{def:normalized-collection}
A collection \(\mathcal S\) of Seregin ancient limits is called normalized if
there are a compact cylinder \(K\Subset\mathbb R^3\times\mathbb R\) and a
constant \(\eta_0>0\) such that every \(V\in\mathcal S\), with its associated
centered pressure \(P\), admits a measurable time-dependent pressure gauge
\(a_V(\tau)\) for which
\begin{equation}\label{eq:collection-nonvanishing}
   \int_K\left(|V(y,\tau)|^3+|P(y,\tau)-a_V(\tau)|^{3/2}\right)\,dy\,d\tau
   \ge\eta_0 .
\end{equation}
When \(\mathcal S\) is obtained from one finite-energy Type I source, \(K\) and
\(\eta_0\) are fixed by the chosen concentration normalization; the admissible
gauge is allowed to depend on the element \(V\).
The condition is independent of the representative of the pressure class:
replacing \(P\) by \(P-c(\tau)\) only replaces the gauge \(a_V(\tau)\) by
\(a_V(\tau)-c(\tau)\).
The integral is finite on \(K\) for every Seregin limit because local
\(L^{3/2}\)-pressure compactness after admissible gauges is part of
\cite{Paper1TypeI}, lemma \texttt{lem:pressure}, and is preserved by the
centered change of variables in \cite{Paper1TypeI}, lemma
\texttt{lem:renorm-bounds}.
\end{definition}

\begin{definition}[Associated normalized collection]\label{def:associated-collection}
Fix a finite-energy Type I source \((u,p,z_*)\).  An associated normalized
collection is obtained as follows.  First choose one positive local
concentration normalization furnished by
\cite{Paper1TypeI}, lemma \texttt{lem:scale-selection}; equivalently, fix the
compact cylinder and lower bound which persist in the ancient limit by
\cite{Paper1TypeI}, proposition \texttt{prop:nonzero}.  Then let
\(\mathcal S\) be any collection of Seregin ancient limits produced from the
same source by subsequences compatible with that fixed normalization, each
equipped with an admissible centered pressure gauge.  The common cylinder
\(K\) and lower bound \(\eta_0\) are those fixed in the first step.  With these
data, \(\mathcal S\) satisfies \Cref{def:normalized-collection}.  The gauges
are measurable: before centering they are spatial averages over fixed balls,
as in \cite{Paper1TypeI}, lemma \texttt{lem:pressure}, and after centering
they are composed with the smooth time map \(t=-e^{-\tau}\) and multiplied by
the smooth factor \(e^{-\tau}\).
\end{definition}

\begin{remark}\label{rem:normalization}
The nonvanishing conditions \eqref{eq:nonvanishing} and
\eqref{eq:collection-nonvanishing} are not additional assumptions.  They are
the persistence, under the compactness used in the Seregin blow-up procedure,
of the positive CKN concentration supplied by \Cref{thm:local-concentration}.
This persistence is precisely the scale-selection and nonzero-limit mechanism
proved in \cite{Paper1TypeI}, lemma \texttt{lem:scale-selection} and
proposition \texttt{prop:nonzero}.  The sole choice made in forming
\(\mathcal S\) is which concentration normalization is fixed once and for all.
\end{remark}

\begin{lemma}[Centered equation]\label{lem:centered-equation}
Let \((U,\Pi)\) be a smooth ancient Navier--Stokes solution on
\(\mathbb R^3\times(-\infty,0)\), and define
\[
   t=-e^{-\tau},\qquad x=e^{-\tau/2}y,\qquad
   V(y,\tau)=e^{-\tau/2}U(e^{-\tau/2}y,-e^{-\tau}),
\]
with
\[
   P(y,\tau)=e^{-\tau}\Pi(e^{-\tau/2}y,-e^{-\tau}).
\]
Then \(V\) satisfies \eqref{eq:centered} and \(\nabla_y\cdot V=0\).
\end{lemma}

\begin{proof}
This is the self-similar change-of-variables computation carried out in
\cite{Paper1TypeI}, lemma \texttt{lem:renormalized-equation}.  That lemma
differentiates \(U(x,t)=e^{\tau/2}V(e^{\tau/2}x,\tau)\), records the common
factor \(e^{3\tau/2}\) in \(\partial_tU\), \((U\cdot\nabla)U\),
\(\nabla\Pi\), and \(\Delta U\), and verifies that
\(\nabla_x\cdot U=e^\tau\nabla_y\cdot V\).  Dividing by the common factor gives
\eqref{eq:centered}.
\end{proof}

\begin{lemma}[Zero velocity and pressure normalization]\label{lem:zero-pressure}
Let \(V\equiv0\) solve \eqref{eq:centered} in distributions on
\(\mathbb R^3\times\mathbb R\).  Then every associated pressure satisfies
\(\nabla_yP=0\).  Consequently, after subtracting the admissible gauge
given by the resulting spatial constant \(a(\tau)\), the integrand in
\eqref{eq:nonvanishing}, and likewise the integrand in
\eqref{eq:collection-nonvanishing}, vanishes identically.
\end{lemma}

\begin{proof}
With \(V\equiv0\), equation \eqref{eq:centered} reduces to
\(\nabla_yP=0\) in distributions, meaning
\[
   \int_{\mathbb R}\int_{\mathbb R^3}P(y,\tau)\,\nabla_y\cdot\varphi(y,\tau)
      \,dy\,d\tau=0
\]
for every \(\varphi\in C_c^\infty(\mathbb R^3\times\mathbb R;\mathbb R^3)\).
By applying this identity to test fields of the form
\(\varphi(y,\tau)=\psi(\tau)\phi(y)e_j\), with
\(\psi\in C_c^\infty(\mathbb R)\) and
\(\phi\in C_c^\infty(\mathbb R^3)\), one obtains
\[
   \int_{\mathbb R}\psi(\tau)
      \left\langle \partial_{y_j}P(\cdot,\tau),\phi\right\rangle
      \,d\tau=0 .
\]
Since this holds for every \(\psi\), Fubini's theorem and the fundamental
lemma of distributions imply that
\(\partial_{y_j}P(\cdot,\tau)=0\) in \(\mathcal D'(\mathbb R^3)\) for almost
every \(\tau\).  Intersecting the three full-measure sets gives
\(\nabla_yP(\cdot,\tau)=0\) for almost every \(\tau\).  The elementary
distributional fact that a distribution on \(\mathbb R^3\) with zero first
spatial derivatives is constant gives
\[
   P(\cdot,\tau)=a(\tau)\quad\hbox{in }\mathcal D'(\mathbb R^3)
\]
for almost every \(\tau\).  The function \(a\) is measurable because it may be
chosen, for a fixed \(\chi\in C_c^\infty(\mathbb R^3)\) with
\(\int\chi=1\), as
\[
   a(\tau)=\langle P(\cdot,\tau),\chi\rangle .
\]
This is an admissible pressure gauge.  After subtracting it,
\(P-a(\tau)=0\) almost everywhere, while \(|V|^3=0\).  Hence every
nonvanishing integral of the form \eqref{eq:nonvanishing} or
\eqref{eq:collection-nonvanishing} is zero for the zero velocity field.
\end{proof}

\begin{lemma}[Basic properties]\label{lem:basic-properties}
Every Seregin ancient limit \(V\) satisfies
\[
   \|V\|_{L^\infty(\mathbb R^3\times\mathbb R)}<\infty
\]
and is smooth on \(\mathbb R^3\times\mathbb R\).  After fixing the whole-space
Leray pressure gauge, \(V\) is a bounded mild ancient solution of
\eqref{eq:centered}.  Its time translates
\[
   V^s(y,\tau)=V(y,\tau+s),\qquad s\in\mathbb R,
\]
are again bounded mild ancient solutions of \eqref{eq:centered} with the same
\(L^\infty\) norm.
\end{lemma}

\begin{proof}
The \(L^\infty\) bound is the Type I ancient bound written in centered
variables:
\[
   |V(y,\tau)|
   =
   \sqrt{-t}\,|U(y\sqrt{-t},t)|
   \le
   \sup_{s<0}\sqrt{-s}\,\|U(s)\|_{L^\infty}.
\]
Smoothness in centered variables is \cite{Paper1TypeI}, lemma
\texttt{lem:smoothness}.  The whole-space Leray pressure gauge and projected
Duhamel formulation are those of the bounded mild framework in
\cite{Paper3Minimal}: definition \texttt{def:mild} fixes the mild formulation
for \eqref{eq:centered}, lemma \texttt{lem:smoothing} gives the smooth
representative, and the time-translation invariance following that definition
shows that \(V^s(y,\tau)=V(y,\tau+s)\) is again a bounded mild ancient solution
with the same \(L^\infty\) norm.  The centered equation is autonomous, so no
additional normalization is introduced by translating \(\tau\).
\end{proof}

\begin{lemma}[Local energy bounds in centered variables]\label{lem:finite-energy-local}
Let \(V\) be a Seregin ancient limit.  For every \(R<\infty\) and every compact
interval \(I\subset\mathbb R\), one has
\[
   \int_I\int_{B_R}|V(y,\tau)|^2\,dy\,d\tau
   +
   \int_I\int_{B_R}|\nabla V(y,\tau)|^2\,dy\,d\tau
   <\infty .
\]
If \(P\) is a pressure associated with \(V\), then after subtracting a function
of \(\tau\) one also has
\[
   P\in L^{3/2}(B_R\times I).
\]
These local bounds are stable under locally smooth limits of time translates
on each fixed compact cylinder.
\end{lemma}

\begin{proof}
The unrenormalized ancient limit \((U,\Pi)\) is suitable and satisfies the
local energy and pressure bounds on compact negative-time cylinders by
\cite{Paper1TypeI}, lemmas \texttt{lem:energy-diss} and \texttt{lem:pressure}
and proposition \texttt{prop:ancient-limit}.  The transfer to centered
variables is the bounded-renormalization estimate of \cite{Paper1TypeI}, lemma
\texttt{lem:renorm-bounds}: on every compact \((y,\tau)\)-cylinder the map
\[
   (y,\tau)\mapsto (x,t)=(e^{-\tau/2}y,-e^{-\tau})
\]
has Jacobian \(dx\,dt=e^{-5\tau/2}\,dy\,d\tau\), with
\(V=e^{-\tau/2}U\), \(\nabla_yV=e^{-\tau}\nabla_xU\), and
\(P=e^{-\tau}\Pi\).  Since the exponential weights are bounded above and below
on compact \(\tau\)-intervals, the local \(L^2\), dissipation, and
\(L^{3/2}\)-pressure estimates pass in the stated form.  Stability under
locally smooth velocity limits and weak local pressure limits is the following
consequence of the same compactness package.  If
\(V_n\to V\) in \(C^\infty_{\mathrm{loc}}\) and the associated pressures,
after admissible gauges, converge weakly in
\(L^{3/2}_{\mathrm{loc}}\), then the \(L^2\) and \(L^2\)-gradient bounds pass
on each \(B_R\times I\) by ordinary convergence of the integrands
\(|V_n|^2\) and \(|\nabla V_n|^2\), while the pressure bound passes by the weak
lower semicontinuity of the \(L^{3/2}(B_R\times I)\)-norm.
\end{proof}

\begin{remark}\label{rem:no-global-tightness}
\Cref{lem:finite-energy-local} is deliberately local.  It does not give
uniform \(L^3\)-tightness in \(y\), nor does it give a global
\(L^\infty_\tau L^3_y\) bound.  The absence of this information is the reason
for the residual class below.
\end{remark}

\section{Admissible classes and the residual class}

Let \(\mathcal S\) denote a normalized collection of Seregin ancient limits in
the sense of \Cref{def:normalized-collection}.  Thus all elements of
\(\mathcal S\) are associated with the same fixed nonvanishing normalization
\((K,\eta_0)\), while their pressure gauges may differ.  Every class introduced
in this section is a subset of this fixed ambient set \(\mathcal S\), and every
set-theoretic complement is taken relative to \(\mathcal S\).  For the
remainder of the paper, fix admissible parameter sets
\[
   I,J\subset(0,\infty),
\]
where admissible means nonempty.

If \(V\in\mathcal S\) is represented by an unrenormalized ancient solution
\(U\), then changing the extraction only within the same Seregin limit changes
neither the velocity \(V\) nor the physical velocity \(U\) determined by
\[
   U(x,t)=(-t)^{-1/2}V(x/\sqrt{-t},-\log(-t)).
\]
Thus the structural properties below are properties of the velocity field, not
of the pressure gauge or of a particular pressure representative.
All pointwise-in-time conditions in the definitions below are imposed on the
smooth representative supplied by \Cref{lem:basic-properties}.  Hence
statements such as \(V(y,\tau)=W(y)\), \(L^\infty_\tau L^3_y\)-bounds, and
uniform tail limits are literal statements for the smooth ancient solution,
not merely assertions about a choice of time slices modulo a null set.

\begin{definition}[Small-amplitude class]\label{def:small-branch}
Let \(\varepsilon_0>0\) be the absolute constant in the small bounded ancient
Liouville theorem for \eqref{eq:centered}.  An element \(V\in\mathcal S\)
belongs to \(\mathcal X_{\mathrm{small}}\) if
\[
   \|V\|_{L^\infty(\mathbb R^3\times\mathbb R)}\le\varepsilon_0 .
\]
\end{definition}

\begin{definition}[Stationary \(L^3\) class]\label{def:stationary-branch}
An element \(V\in\mathcal S\) belongs to
\(\mathcal X_{\mathrm{stat}\text{-}L^3}\) if it is stationary in centered
variables,
\[
   V(y,\tau)=W(y),
\]
and
\[
   W\in L^3(\mathbb R^3).
\]
\end{definition}

\begin{definition}[Uniformly tight class]\label{def:tight-branch}
An element \(V\in\mathcal S\) belongs to
\(\mathcal X_{L^3\mathrm{-tight}}\) if
\begin{equation}\label{eq:L3-bound}
   \sup_{\tau\in\mathbb R}\|V(\tau)\|_{L^3(\mathbb R^3)}<\infty
\end{equation}
and
\begin{equation}\label{eq:L3-tightness}
   \lim_{R\to\infty}
   \sup_{\tau\in\mathbb R}\int_{|y|>R}|V(y,\tau)|^3\,dy=0 .
\end{equation}
\end{definition}

\begin{definition}[Fast-decay class]\label{def:fast-branch}
Fix an admissible set \(J\subset(0,\infty)\).  An element \(V\in\mathcal S\)
belongs to \(\mathcal X_{\mathrm{fast}}(J)\) if, for some \(m\in J\),
\begin{equation}\label{eq:weighted-decay}
   \sup_{\tau\in\mathbb R}
   \int_{\mathbb R^3}(1+|y|^2)^m |V(y,\tau)|^2\,dy<\infty .
\end{equation}
\end{definition}

\begin{definition}[Structured class]\label{def:structured-branch}
Fix admissible parameter sets \(I,J\subset(0,\infty)\).  An element
\(V\in\mathcal S\) belongs to \(\mathcal X_{\mathrm{str}}(I,J)\) if its
physical ancient representative \(U\) satisfies at least one of the following:
\begin{enumerate}
\item \(U\) is axisymmetric without swirl;
\item \(U\) is axisymmetric and, for some \(\gamma\in I\),
\[
   \sup_{\tau\in\mathbb R}\sup_{\rho>0,\ Z\in\mathbb R}
   \rho^\gamma |V_\theta(\rho,Z,\tau)|<\infty;
\]
\item \(U\) is axisymmetric and, for some \(1\le q<\infty\),
\[
   rU_\theta\in L_t^\infty L_x^q
      \bigl(\mathbb R^3\times(-\infty,0)\bigr);
\]
\item \(U\) is axisymmetric, periodic in the axial variable \(z\), and
\[
   rU_\theta\in L^\infty\bigl(\mathbb R^3\times(-\infty,0)\bigr);
\]
\item \(V\) satisfies the weighted decay condition \eqref{eq:weighted-decay}
for some \(m\in J\).
\end{enumerate}
\end{definition}

\begin{remark}\label{rem:structured-import}
The structured alternatives are the centered-limit formulation of the
axisymmetric and weighted-decay branches treated in \cite{Paper6Structured}.
The preservation of the stated hypotheses under self-similar variables and
time shifts is \cite{Paper6Structured}, lemma
\texttt{lem:structure-preserved}.  The no-swirl, \(L_t^\infty L_x^p\)-swirl,
periodic bounded-\(\Gamma\), controlled-swirl, and weighted-decay exclusions
are recorded there as propositions \texttt{prop:no-swirl-type-i},
\texttt{prop:lp-swirl-exclusion}, \texttt{prop:periodic-swirl-exclusion},
\texttt{prop:swirl-exclusion}, and \texttt{prop:weighted-exclusion}.  The
present paper uses those mechanisms solely through the single structured
Liouville input stated below.
\end{remark}

\begin{definition}[Residual Seregin class]\label{def:residual}
The residual class associated with \(\mathcal S\), \(I\), and \(J\) is
\[
   \mathcal R(\mathcal S;I,J)
   =
   \mathcal S
   \setminus
   \left(
      \mathcal X_{\mathrm{small}}
      \cup
      \mathcal X_{\mathrm{stat}\text{-}L^3}
      \cup
      \mathcal X_{L^3\mathrm{-tight}}
      \cup
      \mathcal X_{\mathrm{fast}}(J)
      \cup
      \mathcal X_{\mathrm{str}}(I,J)
   \right).
\]
Equivalently, \(V\in\mathcal R(\mathcal S;I,J)\) is a Seregin ancient limit
which is neither small in \(L^\infty\), nor stationary and globally \(L^3\),
nor uniformly \(L^3\)-tight, nor fast decaying in an admissible weighted
class, nor contained in the stated structured class.
\end{definition}

\begin{remark}\label{rem:overlap}
The classes are not intended to be disjoint.  For example, a weighted decay
condition may imply an \(L^3\)-tightness statement under additional boundedness
and interpolation.  The residual class is defined as the complement of the
union, so overlaps among the non-residual classes cause no ambiguity.
\end{remark}

\begin{proposition}[Exhaustion by admissible classes]\label{prop:exhaustion}
Let \(\mathcal S\) be a normalized collection of Seregin ancient limits, and
fix admissible sets \(I,J\subset(0,\infty)\).  Every \(V\in\mathcal S\)
satisfies at least one of the following alternatives:
\[
\begin{gathered}
   V\in\mathcal X_{\mathrm{small}},\qquad
   V\in\mathcal X_{\mathrm{stat}\text{-}L^3},\qquad
   V\in\mathcal X_{L^3\mathrm{-tight}},\\
   V\in\mathcal X_{\mathrm{fast}}(J),\qquad
   V\in\mathcal X_{\mathrm{str}}(I,J),\qquad
   V\in\mathcal R(\mathcal S;I,J).
\end{gathered}
\]
\end{proposition}

\begin{proof}
Set
\[
   \mathcal U
   =
   \mathcal X_{\mathrm{small}}
   \cup
   \mathcal X_{\mathrm{stat}\text{-}L^3}
   \cup
   \mathcal X_{L^3\mathrm{-tight}}
   \cup
   \mathcal X_{\mathrm{fast}}(J)
   \cup
   \mathcal X_{\mathrm{str}}(I,J).
\]
For any \(V\in\mathcal S\), either \(V\in\mathcal U\), in which case \(V\)
belongs to at least one of the five displayed classes by the definition of
union, or \(V\notin\mathcal U\).  In the second case
\[
   V\in\mathcal S\setminus\mathcal U
   =
   \mathcal R(\mathcal S;I,J)
\]
by \Cref{def:residual}.  This proves the stated exhaustive alternative.
\end{proof}

\section{Liouville inputs for the non-residual classes}

We now state the Liouville inputs used in the reduction.  The first two
closures are unconditional: small bounded ancient solutions vanish by the
centered mild formulation, and stationary \(L^3\) self-similar profiles vanish
by the theorem of Necas--Ruzicka--Sverak.  The remaining inputs are hypotheses
of the present theorem, formulated so that each can be replaced by any
independently proved Liouville theorem in the corresponding class.

\begin{lemma}[Small bounded ancient Liouville theorem]
\label{lem:small-liouville}
There is an absolute \(\varepsilon_0>0\) such that every bounded mild ancient
solution \(V\) of \eqref{eq:centered} with
\[
   \|V\|_{L^\infty(\mathbb R^3\times\mathbb R)}\le\varepsilon_0
\]
is identically zero.
\end{lemma}

\begin{proof}
Write the projected centered equation as
\[
   \partial_\tau V=LV-\mathbb P\nabla\cdot(V\otimes V),
   \qquad
   L=\Delta-\frac12y\cdot\nabla-\frac12 .
\]
Let \(S(s)=e^{sL}\) be the centered Stokes semigroup on divergence-free
fields.  The Ornstein--Uhlenbeck formula for \(S(s)\) and its \(L^\infty\)
decay estimate are \cite{Paper3Minimal}, equations \texttt{eq:OU-semigroup}
and \texttt{eq:stokes-Linfty-estimate}:
\[
   (S(s)f)(y)
   =
   e^{-s/2}
   \left(e^{(1-e^{-s})\Delta}f\right)(e^{-s/2}y),
   \qquad s>0.
\]
Hence
\[
   \|S(s)f\|_{L^\infty}\le e^{-s/2}\|f\|_{L^\infty}.
\]
The corresponding projected-divergence estimate is the bounded-data Stokes
estimate recorded in \cite{Paper3Minimal}, equation
\texttt{eq:stokes-\allowbreak Linfty-\allowbreak estimate}; the adjoint-test
estimate \texttt{eq:stokes-\allowbreak adjoint-\allowbreak test-\allowbreak estimate}
in the same paper is the dual formulation used to define the weak-* Duhamel
term.  These are the centered forms of Kato's \(L^p\) estimates
\cite{Kato1984}.  Thus there is an absolute constant \(C_0\) such that, for
every bounded tensor field \(F\),
\[
   \|S(s)\mathbb P\nabla\cdot F\|_{L^\infty}
   \le
   C_0 e^{-s/2}(1-e^{-s})^{-1/2}\|F\|_{L^\infty},
   \qquad s>0.
\]
The factor on the right is integrable on \((0,\infty)\).  Set
\[
   A=C_0\int_0^\infty e^{-s/2}(1-e^{-s})^{-1/2}\,ds<\infty .
\]
Indeed, \(1-e^{-s}\sim s\) as \(s\downarrow0\), so the kernel is comparable to
\(s^{-1/2}\) near the origin, while the factor \(e^{-s/2}\) gives exponential
decay as \(s\to\infty\).

Fix \(\tau\in\mathbb R\) and \(T>0\).  The mild identity is definition
\texttt{def:mild} of \cite{Paper3Minimal}.  Applied on \([\tau-T,\tau]\), it
gives
\[
   V(\tau)
   =
   S(T)V(\tau-T)
   -
   \int_0^T
      S(T-\sigma)\mathbb P\nabla\cdot
      \bigl(V\otimes V\bigr)(\tau-T+\sigma)\,d\sigma .
\]
If \(M=\|V\|_{L^\infty(\mathbb R^3\times\mathbb R)}\), the preceding two
estimates and the pointwise bound
\(\|V(\tau)\otimes V(\tau)\|_{L^\infty}\le M^2\) imply
\[
\begin{aligned}
   \|V(\tau)\|_{L^\infty}
   &\le
   \|S(T)V(\tau-T)\|_{L^\infty}  \\
   &\quad+
   \int_0^T
      \left\|
      S(T-\sigma)\mathbb P\nabla\cdot
      \bigl(V\otimes V\bigr)(\tau-T+\sigma)
      \right\|_{L^\infty}\,d\sigma        \\
   &\le
   e^{-T/2}M
   +
   C_0 M^2\int_0^T
      e^{-(T-\sigma)/2}
      \bigl(1-e^{-(T-\sigma)}\bigr)^{-1/2}\,d\sigma .
\end{aligned}
\]
The change of variables \(s=T-\sigma\) and the definition of \(A\) give
\[
   \|V(\tau)\|_{L^\infty}
   \le
   e^{-T/2}M
   +
   A M^2 .
\]
Letting \(T\to\infty\) yields \(\|V(\tau)\|_{L^\infty}\le A M^2\).  Taking
the supremum in \(\tau\) gives \(M\le A M^2\).  If \(M>0\), division by \(M\)
gives \(1\le AM\).  Choose \(\varepsilon_0<1/A\).  The assumption
\(M\le\varepsilon_0\) then contradicts \(1\le AM\).  Hence \(M=0\), which is
equivalent to \(V\equiv0\).
\end{proof}

\begin{theorem}[Unconditional small and stationary branch exclusion]
\label{thm:classical-closed-branches}
Let \(\mathcal S\) be a normalized collection of Seregin ancient limits, and
let \(V\in\mathcal S\) be nonzero.  Then
\[
   V\notin\mathcal X_{\mathrm{small}}
   \qquad\hbox{and}\qquad
   V\notin\mathcal X_{\mathrm{stat}\text{-}L^3}.
\]
\end{theorem}

\begin{proof}
Assume first that \(V\in\mathcal X_{\mathrm{small}}\).  By definition,
\[
   \|V\|_{L^\infty(\mathbb R^3\times\mathbb R)}\le\varepsilon_0 ,
\]
where \(\varepsilon_0\) is the small-data constant for bounded mild ancient
solutions of the centered equation from \Cref{lem:small-liouville}.  The mild
solution property is part of \Cref{lem:basic-properties}, hence that lemma
gives \(V\equiv0\).  By \Cref{lem:zero-pressure}, after subtracting
the allowed pressure gauge the integral in \eqref{eq:nonvanishing} is zero.
This contradicts the inherited positive normalization.

Assume next that \(V\in\mathcal X_{\mathrm{stat}\text{-}L^3}\).  Then
\(V(y,\tau)=W(y)\) for some \(W\in L^3(\mathbb R^3)\).  Since \(V\) is a
smooth distributional solution of \eqref{eq:centered}, the identity
\(\partial_\tau V=0\) is classical and the remaining terms in the equation are
smooth functions of \(y\).  For every
\(\phi\in C_c^\infty(\mathbb R^3;\mathbb R^3)\) and
\(\psi\in C_c^\infty(\mathbb R)\), testing the equation against
\(\psi(\tau)\phi(y)\) gives
\[
   \int_{\mathbb R}\psi(\tau)
   \left\langle
      \nabla_yP(\cdot,\tau)
      -\Delta W+\frac12(W+y\cdot\nabla W)+(W\cdot\nabla)W,
      \phi
   \right\rangle\,d\tau=0 .
\]
Taking a countable dense family of spatial tests and using the fundamental
lemma of distributions in \(\tau\), there is a full-measure set
\(E\subset\mathbb R\) such that, for every \(\tau\in E\),
\[
   \nabla_yP(\cdot,\tau)
   =
   \Delta W-\frac12(W+y\cdot\nabla W)-(W\cdot\nabla)W
   \quad\hbox{in }\mathcal D'(\mathbb R^3).
\]
The right-hand side is independent of \(\tau\).  Hence for
\(\tau_1,\tau_2\in E\),
\[
   \nabla_y\bigl(P(\cdot,\tau_1)-P(\cdot,\tau_2)\bigr)=0
   \quad\hbox{in }\mathcal D'(\mathbb R^3),
\]
so \(P(\cdot,\tau_1)-P(\cdot,\tau_2)\) is spatially constant.  Fix
\(\tau_0\in E\), set \(P_{\mathrm{st}}=P(\cdot,\tau_0)\), and subtract the
admissible time-dependent constant
\[
   a(\tau)=\langle P(\cdot,\tau)-P_{\mathrm{st}},\chi\rangle ,
   \qquad \int_{\mathbb R^3}\chi=1,
\]
for a fixed \(\chi\in C_c^\infty(\mathbb R^3)\).  Then the gauged pressure
equals the single stationary pressure profile \(P_{\mathrm{st}}\) for
almost every \(\tau\in E\), and the profile equation is
\[
   (W\cdot\nabla)W+\nabla P_{\mathrm{st}}
   =
   \Delta W-\frac12(W+y\cdot\nabla W),
   \qquad \nabla\cdot W=0 .
\]
The stationary-profile passage and the \(L^3\) Liouville input are recorded in
\cite{Paper5Equilibrium}.  The lemma labelled \texttt{lem:stationary}-\texttt{profile}
identifies stationary mild trajectories with stationary self-similar profiles,
and theorem \texttt{thm:NRS} states the Necas--Ruzicka--Sverak theorem
\cite{NRS1996} in this form.  Thus the
assumed \(W\in L^3(\mathbb R^3)\) gives \(W\equiv0\).  Again
\Cref{lem:zero-pressure} makes the left-hand side of
\eqref{eq:nonvanishing} vanish, contradicting the inherited normalization.
\end{proof}

\begin{hypothesis}[Tight-class Liouville input]\label{hyp:tight-liouville}
For every normalized Seregin collection \(\mathcal S\), no nonzero element of
\(\mathcal S\cap\mathcal X_{L^3\mathrm{-tight}}\) exists.
\end{hypothesis}

\begin{hypothesis}[Structured Liouville input]\label{hyp:structured-liouville}
For every normalized Seregin collection \(\mathcal S\), no nonzero element of
\(\mathcal S\cap\mathcal X_{\mathrm{str}}(I,J)\) exists.
\end{hypothesis}

\begin{remark}\label{rem:hypotheses}
\Cref{hyp:tight-liouville} is the form obtained, for instance, from a
minimal-element reduction plus compact-hull equilibrium rigidity in the
uniformly \(L^3\)-tight class.  \Cref{hyp:structured-liouville} contains the
structured reductions imported in \Cref{rem:structured-import}: the
axisymmetric no-swirl Liouville theorem of Koch--Nadirashvili--Seregin--Sverak
\cite{KNSS2009}, the \(L_t^\infty L_x^p\)-swirl theorem of Lei--Zhang--Zhao
\cite{LZZ2017}, the periodic bounded-\(\Gamma\) theorem of Lei--Ren--Zhang
\cite{LRZ2022}, and any additional admissible controlled-swirl or
weighted-decay Liouville criteria assumed for the sets \(I\) and \(J\).
\end{remark}

\begin{proposition}[Non-residual class exclusion]\label{prop:closed-branch-exclusion}
Assume \Cref{hyp:tight-liouville,hyp:structured-liouville}.  Let
\(\mathcal S\) be a normalized collection of Seregin ancient limits, and let
\(V\in\mathcal S\) be nonzero.  Then \(V\) cannot belong to
\[
   \mathcal X_{\mathrm{small}}
   \cup
   \mathcal X_{\mathrm{stat}\text{-}L^3}
   \cup
   \mathcal X_{L^3\mathrm{-tight}}
   \cup
   \mathcal X_{\mathrm{fast}}(J)
   \cup
   \mathcal X_{\mathrm{str}}(I,J).
\]
\end{proposition}

\begin{proof}
The small-amplitude and stationary \(L^3\) classes are excluded by
\Cref{thm:classical-closed-branches}.
Membership in \(\mathcal X_{L^3\mathrm{-tight}}\) is excluded by
\Cref{hyp:tight-liouville}.  Membership in
\(\mathcal X_{\mathrm{str}}(I,J)\) is excluded by
\Cref{hyp:structured-liouville}.  Finally,
\(\mathcal X_{\mathrm{fast}}(J)\subset\mathcal X_{\mathrm{str}}(I,J)\) because
both definitions require the existence of some \(m\in J\) for which
\eqref{eq:weighted-decay} holds, and this is precisely the fifth alternative in
\Cref{def:structured-branch}.  Thus the structured hypothesis also excludes
the fast-decay class.
\end{proof}

\section{The residual class}

The residual class is the remaining alternative not covered by the stated
Liouville inputs.  We state the needed assertion in two equivalent forms.

\begin{hypothesis}[No residual Seregin class]\label{hyp:no-residual}
For the fixed admissible sets \(I,J\), every finite-energy Type I source and
every associated normalized collection \(\mathcal S\) in the sense of
\Cref{def:associated-collection} satisfy
\[
   \mathcal R(\mathcal S;I,J)=\varnothing .
\]
\end{hypothesis}

\begin{hypothesis}[Residual rigidity]\label{hyp:residual-rigidity}
For the fixed admissible sets \(I,J\), every finite-energy Type I source,
every associated normalized collection \(\mathcal S\) in the sense of
\Cref{def:associated-collection}, and every element of
\(\mathcal R(\mathcal S;I,J)\) is identically zero.
\end{hypothesis}

\begin{proposition}[Equivalence for Seregin limits]\label{prop:residual-equivalence}
For normalized Seregin collections associated with finite-energy Type I
sources in the sense of \Cref{def:associated-collection}, the no-residual
assertion \Cref{hyp:no-residual} and the residual rigidity assertion
\Cref{hyp:residual-rigidity} are equivalent.
\end{proposition}

\begin{proof}
If \Cref{hyp:no-residual} holds, then the residual class is empty and the
rigidity assertion is vacuous.  Conversely, assume
\Cref{hyp:residual-rigidity}, and let \(\mathcal S\) be a normalized Seregin
collection associated with a finite-energy Type I source in the sense of
\Cref{def:associated-collection}.  If
\(\mathcal R(\mathcal S;I,J)=\varnothing\), there is nothing to prove.  If it
is nonempty, choose \(V\in\mathcal R(\mathcal S;I,J)\).  Then
\Cref{hyp:residual-rigidity} gives \(V\equiv0\).  By \Cref{lem:zero-pressure},
the associated pressure is spatially constant after gauge normalization, so
the left-hand side of \eqref{eq:collection-nonvanishing} is zero.  This
contradicts the defining lower bound
\(\eta_0>0\) of the normalized collection.  Hence
\(\mathcal R(\mathcal S;I,J)=\varnothing\), which is
\Cref{hyp:no-residual}.
\end{proof}

\begin{remark}\label{rem:meaning-residual}
Analytically, \Cref{hyp:no-residual} may be proved in more than one way.  One
route is to show that finite energy and suitability force every Seregin limit
to be uniformly \(L^3\)-tight or fast decaying in self-similar variables.
Another is to prove a Liouville theorem which applies directly to bounded
ancient limits without tightness.  Both routes are represented by the same
formal statement above because either one eliminates
\(\mathcal R(\mathcal S;I,J)\).
\end{remark}

\section{Conditional Type I exclusion}

\begin{theorem}[Residual-class criterion for Type I exclusion]\label{thm:full-type-i}
For the fixed admissible sets \(I,J\), assume the tight-class Liouville input
\Cref{hyp:tight-liouville}, the structured Liouville input
\Cref{hyp:structured-liouville}, and the no-residual assertion
\Cref{hyp:no-residual}.  Then no finite-energy suitable weak solution of the
three-dimensional Navier--Stokes equations develops a Type I singularity in the
sense of \Cref{def:source}.
\end{theorem}

\begin{proof}
Suppose, toward a contradiction, that \((u,p)\) has a finite-energy Type I
singular point \(z_*=(x_*,T)\).  By \Cref{prop:seregin-extraction}, there
exists a Seregin ancient limit \(V\) associated with \(z_*\), satisfying the
inherited nonvanishing condition \eqref{eq:nonvanishing}.  This condition
implies \(V\not\equiv0\): if \(V\equiv0\), then
\Cref{lem:zero-pressure} would make the left-hand side of
\eqref{eq:nonvanishing} vanish.  Fix the concentration normalization used to
obtain this limit, and let \(\mathcal S\) be the collection of all Seregin
limits obtained from the same source with that fixed normalization.  By
\Cref{def:associated-collection,def:normalized-collection} and
\Cref{rem:normalization}, this \(\mathcal S\) is an associated normalized
collection; in particular it is nonempty because it contains the limit \(V\).

By \Cref{prop:exhaustion}, the following exhaustive assertion holds: either
\(V\) belongs to at least one of the five non-residual classes, or
\[
   V\in\mathcal R(\mathcal S;I,J).
\]
If \(V\) belongs to a non-residual class, then
\Cref{prop:closed-branch-exclusion} excludes it.  This proposition invokes the
unconditional small and stationary exclusions for the first two classes,
\Cref{hyp:tight-liouville} for the uniformly \(L^3\)-tight class, and
\Cref{hyp:structured-liouville} for the structured and fast-decay classes.
If \(V\) belongs to the residual class, then
\Cref{hyp:no-residual} asserts that no such element exists.  Both alternatives
are impossible.  This contradiction proves that the assumed Type I singular
point cannot exist.
\end{proof}

\begin{corollary}[Replacement-rigidity version]\label{cor:replacement}
For the fixed admissible sets \(I,J\), the conclusion of
\Cref{thm:full-type-i} remains valid if \Cref{hyp:no-residual} is replaced by
\Cref{hyp:residual-rigidity}.
\end{corollary}

\begin{proof}
Under \Cref{hyp:residual-rigidity}, \Cref{prop:residual-equivalence} gives
\Cref{hyp:no-residual} for the same fixed admissible sets \(I,J\) and the same
associated normalized Seregin collections.  The hypotheses of
\Cref{thm:full-type-i} are therefore satisfied, and the theorem gives the
claimed exclusion of Type I singularities.
\end{proof}

\section{Conclusion}

The criterion isolates the obstruction to finite-energy Type I exclusion as a
single residual statement.  The blow-up construction produces a nonzero
Seregin ancient limit.  The exhaustion result places this limit either in a
class eliminated by the unconditional and assumed Liouville inputs or in the
residual class.
Thus the remaining analytic problem is to prove that finite-energy Seregin
limits cannot persist in the residual non-tight/non-decaying regime, or to
prove a Liouville theorem which applies to that regime directly.

\begin{thebibliography}{99}

\bibitem{AlbrittonBarker2019}
D. Albritton and T. Barker,
\emph{On local Type I singularities of the Navier--Stokes equations and
Liouville theorems},
J. Math. Fluid Mech. \textbf{21} (2019), no. 3, Paper No. 43.

\bibitem{CKN1982}
L. Caffarelli, R. Kohn, and L. Nirenberg,
\emph{Partial regularity of suitable weak solutions of the Navier--Stokes equations},
Comm. Pure Appl. Math. \textbf{35} (1982), no. 6, 771--831.

\bibitem{ESS2003}
L. Escauriaza, G. Seregin, and V. \v{S}verak,
\emph{\(L_{3,\infty}\)-solutions of Navier--Stokes equations and backward uniqueness},
Russian Math. Surveys \textbf{58} (2003), no. 2, 211--250.

\bibitem{KNSS2009}
G. Koch, N. Nadirashvili, G. Seregin, and V. Sverak,
\emph{Liouville theorems for the Navier-Stokes equations and applications},
Acta Math. \textbf{203} (2009), no. 1, 83--105.

\bibitem{Kato1984}
T. Kato,
\emph{Strong \(L^p\)-solutions of the Navier--Stokes equation in
\(\mathbb R^m\), with applications to weak solutions},
Math. Z. \textbf{187} (1984), no. 4, 471--480.

\bibitem{LRZ2022}
Z. Lei, X. Ren, and Q. S. Zhang,
\emph{A Liouville theorem for axi-symmetric Navier--Stokes equations on
\(\mathbb R^2\times\mathbb T^1\)},
Math. Ann. \textbf{383} (2022), no. 1--2, 415--431.

\bibitem{LZZ2017}
Z. Lei, Q. Zhang, and N. Zhao,
\emph{Improved Liouville theorems for axially symmetric Navier--Stokes
equations},
Sci. Sin. Math. \textbf{47} (2017), no. 10, 1183--1198.

\bibitem{NRS1996}
J.~Ne\v{c}as, M.~R\r{u}\v{z}i\v{c}ka, and V.~\v{S}ver\'ak,
\emph{On Leray's self-similar solutions of the Navier--Stokes equations},
Acta Math. \textbf{176} (1996), 283--294.

\bibitem{Paper0LocalConcentration}
Guillem Cordero,
\emph{Local CKN Concentration and the Type I/Type II Dichotomy for
Navier--Stokes},
first-nodes Paper 0,
\texttt{docs/source/first\_nodes/paper0\_local\_concentration\_entry\_typeII.tex}.

\bibitem{Paper1TypeI}
Type I Paper 1,
\emph{Type I Blow-up Limits for the Three-dimensional Navier--Stokes Equations
and the Centered Ancient Equation},
\texttt{docs/source/type\_1/paper1\_typeI\_blowup\_limits\_centered\_ancient.tex}.

\bibitem{Paper3Minimal}
G. Duran Ballester,
\emph{Minimal Ancient Solutions for the Navier--Stokes Equations in
Self-similar Variables},
Type I Paper 3,
\texttt{docs/source/type\_1/paper3\_minimal\_ancient\_solutions\_reduction.tex}.

\bibitem{Paper5Equilibrium}
G. Duran Ballester,
\emph{A Conditional Equilibrium Criterion for Compact Ancient Navier--Stokes
Dynamics},
Type I Paper 5,
\texttt{docs/source/type\_1/paper5\_equilibrium\_rigidity\_compact\_ancient\_dynamics.tex}.

\bibitem{Paper6Structured}
G. Duran Ballester,
\emph{Ancient Navier--Stokes solutions under axisymmetry and decay},
Type I Paper 6,
\texttt{docs/source/type\_1/paper6\_structured\_bounded\_ancient\_solutions.tex}.

\bibitem{Seregin2012}
G. Seregin,
\emph{A certain necessary condition of potential blow up for Navier--Stokes equations},
Comm. Math. Phys. \textbf{312} (2012), no. 3, 833--845.

\end{thebibliography}

\end{document}
